% Two-Aspect Monism Diagram
% Single object Γ with two projections: Physics and Experience

\begin{tikzpicture}[
  box/.style={rectangle, rounded corners=4pt, draw=#1, fill=#1!10,
              minimum width=3.5cm, minimum height=1.2cm, align=center, font=\small},
  gammabox/.style={rectangle, rounded corners=6pt, draw=plospurple, fill=plospurple!15,
                   minimum width=4cm, minimum height=1.5cm, align=center, font=\large\bfseries},
  arrow/.style={-{Stealth[length=3mm]}, thick, #1},
  scale=1.0
]

% === Central object ===
\node[gammabox] (gamma) at (0,0) {$\Gamma \in \mathcal{D}(\mathbb{C}^7)$\\[2pt]
  \footnotesize\color{plosgray} The single primitive};

% === External projection (left) ===
\node[box=plosblue] (physics) at (-4.5, -3) {\textbf{External aspect}\\Physics};
\node[font=\footnotesize, text width=3.8cm, align=center] at (-4.5, -4.5) {
  Eigenvalues, dynamics\\
  Quantum mechanics\\
  Spacetime geometry\\
  Einstein equations
};

% === Internal projection (right) ===
\node[box=red!60!black] (experience) at (4.5, -3) {\textbf{Internal aspect}\\Experience};
\node[font=\footnotesize, text width=3.8cm, align=center] at (4.5, -4.5) {
  Qualia, feelings\\
  Spectral decomp. of $\rho_E$\\
  Fubini--Study metric\\
  The sense of being
};

% === Arrows ===
\draw[arrow=plosblue] (gamma.south west) -- (physics.north)
  node[midway, left, font=\footnotesize] {$\mathrm{Map}_{\mathrm{ext}}$};
\draw[arrow=red!60!black] (gamma.south east) -- (experience.north)
  node[midway, right, font=\footnotesize] {$\mathrm{Map}_{\mathrm{int}}$};

% === NOT a bridge ===
\draw[dashed, gray, thick] (physics.east) -- (experience.west);
\node[font=\footnotesize\color{gray}, align=center] at (0, -3.4)
  {no bridge needed --- \textit{same object}};

% === Philosophical lineage ===
\node[font=\scriptsize, text width=8cm, align=center, color=plosgray] at (0, 2) {
  Spinoza (1677): ``The order of ideas = the order of things''\\
  Russell (1927): ``Neutral stuff, neither mental nor physical''\\
  \textbf{UHM (2026): $\Gamma$ is both --- one object, two projections}
};

\end{tikzpicture}
